\documentclass[11pt]{article}
\usepackage{preamble}
\setlength{\parskip}{4pt}

\begin{document}

\daybanner{AP Statistics \textperiodcentered\ Unit 3 \textperiodcentered\ Topic 3.4 \textperiodcentered\ B: 2026-12-01 \textperiodcentered\ E: 2027-01-06 \textperiodcentered\ TEACHER KEY}{One Interval, Two Threshold Claims}

\sectionbanner{CED TETHER}

{\small
\begin{itemize}[leftmargin=1.5em,itemsep=2pt,topsep=2pt]
  \item Enduring Understanding: UNC-4 An interval of values should be used to estimate parameters, in order to account for uncertainty.
  \item Skills:
  \item \textbf{Skill 4.B} --- Interpret statistical calculations and findings to assign meaning or assess a claim.
  \item \textbf{Skill 4.D} --- Justify a claim based on a confidence interval.
  \item \textbf{Skill 4.A} --- Make an appropriate claim or draw an appropriate conclusion.
  \item \textbf{LO UNC-4.F} --- Interpret a confidence interval for a population proportion. [Skill 4.B]
  \item \textbf{EK UNC-4.F.1} --- A confidence interval for a population proportion either contains the population proportion or it does not, because each interval is based on random sample data, which varies from sample to sample.
  \item \textbf{EK UNC-4.F.2} --- We are C\% confident that the confidence interval for a population proportion captures the population proportion.
  \item \textbf{EK UNC-4.F.3} --- In repeated random sampling with the same sample size, approximately C\% of confidence intervals created will capture the population proportion.
  \item \textbf{EK UNC-4.F.4} --- Interpreting a confidence interval for a one-sample proportion should include a reference to the sample taken and details about the population it represents.
  \item \textbf{LO UNC-4.G} --- Justify a claim based on a confidence interval for a population proportion. [Skill 4.D]
  \item \textbf{EK UNC-4.G.1} --- A confidence interval for a population proportion provides an interval of values that may provide sufficient evidence to support a particular claim in context.
  \item \textbf{LO UNC-4.H} --- Identify the relationships between sample size, width of a confidence interval, confidence level, and margin of error for a population proportion. [Skill 4.A]
  \item \textbf{EK UNC-4.H.1} --- When all other things remain the same, the width of the confidence interval for a population proportion tends to decrease as the sample size increases. For a population proportion, the width of the interval is proportional to 1/$\surd$n.
  \item \textbf{EK UNC-4.H.2} --- For a given sample, the width of the confidence interval for a population proportion increases as the confidence level increases.
  \item \textbf{EK UNC-4.H.3} --- The width of a confidence interval for a population proportion is exactly twice the margin of error.
\end{itemize}
}
\textbf{Essential question:} How can one confidence interval support a modest lower-bound claim but not a stronger one?\par
\textbf{DOK-3 skill:} 4.D \textperiodcentered\ \textbf{FRQ pattern:} {\small\texttt{one-\allowbreak{}interval-\allowbreak{}two-\allowbreak{}threshold-\allowbreak{}claims-\allowbreak{}bound-\allowbreak{}the-\allowbreak{}conclusion}}\par
\textbf{DOK rationale:} Part (c) coordinates interval construction, contextual interpretation, and two endpoint comparisons to adjudicate competing conclusions, distinguish plausibility from support, and bound what can be reported.\par

\frameworkphaseheader{Do Now (first take) $\rightarrow$ Explore (video + follow-along) $\rightarrow$ Exit (finish + turn in)}{1 $\rightarrow$ 3}{45}{%
\begin{itemize}[leftmargin=1.5em,itemsep=2pt,topsep=2pt]
  \item Board slide up at the bell. Ask for a point-estimate first take.
  \item Begin the video at minute 5; return at minute 33 to compare the lower endpoint with both thresholds.
  \item Collect at the door. Score part (c) only, E/P/I.
\end{itemize}
}{%
\begin{itemize}[leftmargin=1.5em,itemsep=2pt,topsep=2pt]
  \item Choose which threshold claims have support and cite one number.
  \item Complete the follow-along, finish (a)--(c), revise, and turn in the sheet.
\end{itemize}
}{%
\begin{itemize}[leftmargin=1.5em,itemsep=2pt,topsep=2pt]
  \item What uncertainty disappears if you report only $0.68$?
  \item Which endpoint should you compare with a lower-bound claim?
  \item What does $0.65$ inside the interval establish---and not establish?
\end{itemize}
}{Solo teacher. Circulate ~5 pairs. Require both endpoint comparisons and a bounded conclusion.}

\begin{callout}{\IconWarn\ WATCH FOR}{calloutred}
\begin{itemize}[leftmargin=1.5em,itemsep=2pt,topsep=2pt]
  \item Treating the sample proportion as the population proportion.
  \item Calling an interior threshold a proved lower bound.
  \item Assigning 95 percent probability to the fixed parameter.
\end{itemize}
\end{callout}

\sectionbanner{SCORING PART (c) --- E / P / I}

\textbf{Expected elements}
\begin{itemize}[leftmargin=1.5em,itemsep=2pt,topsep=2pt]
  \item \textbf{judgment-1} (required) --- Supports the more-than-60-percent claim using lower endpoint about 0.627 above 0.60.
  \item \textbf{judgment-2} (required) --- Withholds the at-least-65-percent claim because the lower endpoint is below 0.65, while 0.65 remains plausible.
  \item \textbf{judgment-3} (required) --- Supports Mira and explains that Kai's point estimate alone ignores sampling uncertainty.
\end{itemize}

{\small\renewcommand{\arraystretch}{1.25}
\noindent\begin{tabular}{|p{0.5in}|p{5.9in}|}
\hline
\textbf{E} & Meets all three checks with correct contextual reasoning; equivalent wording is acceptable. \\ \hline
\textbf{P} & Shows at least one substantive correct check, but has an omission or error that prevents E. \\ \hline
\textbf{I} & Shows none of the three checks with substantive correct reasoning; a choice alone is insufficient. \\ \hline
\end{tabular}}

\clearpage
\focusbanner{TODAY'S PROBLEM --- annotated}

Suppose a statewide activities association has 10,000 member students and is considering requiring reusable food containers at its tournaments. In a simple random sample of 300 member students, 204 favor the requirement. A sponsor claims that at least $65\%$ of all member students favor it; an event organizer makes the more modest claim that more than $60\%$ favor it.\par\smallskip \textbf{Kai:} ``The sample has $204/300=68\%$ in favor, above both thresholds, so it supports both claims.''\par \textbf{Mira:} ``The point estimate alone cannot settle both claims. We should construct a confidence interval and compare all of its plausible values with each threshold.''\par
\begin{center}
\textbf{Sample counts}\par\smallskip
\begin{tabular}{|l|c|c|c|c|}
\hline
\textbf{Group} & \textbf{N} & \textbf{n} & \textbf{Favor} & \textbf{Other} \\ \hline
\textbf{Members} & 10,000 & 300 & 204 & 96 \\ \hline
\end{tabular}
\end{center}
\begin{firsttakebox}{FIRST TAKE (5 min)}
Do the data appear to support both claims, only one, or neither? Cite one number and commit before calculating the interval.\par
\answer{Not graded. The interval must separate plausible from justified.}
\end{firsttakebox}

\textit{Rules callout ``LET THE WHOLE INTERVAL SET THE BOUND (keep on the page)'' is printed on the student sheet.}\par\medskip

\dokbadge{1}~\textbf{(a)} Define the population proportion $p$ and calculate the sample proportion $\hat p$.\par
\answer{Let $p$ be the proportion of all 10,000 association member students who favor requiring reusable food containers at tournaments. The sample proportion is $\hat p=204/300=0.68$.}

\dokbadge{2}~\textbf{(b)} Verify the conditions, construct a 95 percent confidence interval for $p$, and interpret the interval in the association's context.\par
\answer{The association used an SRS; $300\le0.10(10{,}000)=1{,}000$; and there are 204 successes and 96 failures, both at least 10. Thus $SE=\sqrt{0.68(0.32)/300}=0.0269320$ and $ME=1.96(SE)=0.0527867$. The 95 percent interval is $(0.6272,0.7328)$, or about $(0.627,0.733)$. We are 95 percent confident that this interval captures the proportion of all association member students who favor the reusable-container requirement.}

\dokbadge{3}~\textbf{(c)} In 3--4 sentences, decide which threshold claim the interval supports and whose reasoning is better. Compare the lower endpoint with 0.60 and 0.65. Explain why a plausible 0.65 is different from establishing at least 0.65.\par
\answer{Mira is warranted. The entire interval $(0.627,0.733)$ exceeds $0.60$, supporting the organizer's claim. Its lower endpoint is below $0.65$, so the interval does not establish the sponsor's lower bound. Because $0.65$ is inside, it remains plausible rather than disproved. Kai treats $0.68$ as exact, hiding the margin of error and falsely suggesting that both population bounds are established.}

\begin{summaryexitbox}{EXIT}
One thing the video changed about my first take: \hrulefill
\end{summaryexitbox}

\end{document}
